\documentstyle[% 


righttag% 


,11pt% 


,amssymb% 


,russian%               remove in Eng. 


,izvru%                 Set izven% in Eng. 


]{amsart} 


 


% 


 


\newcommand{\col}{\operatorname{col}} 


\newcommand{\tts}[1]{} 


\renewcommand{\baselinestretch}{1.05}


 


\theoremstyle{definition} 


\theorembodyfont{\rm} 


\newtheorem{defnnonum}{\hskip\parindent Определение} 


\renewcommand{\thedefnnonum}{} 


 


%\hoffset=-15mm%                  For Eng. 


\setcounter{page}{10} 


\god{2003} 


\nomer{9~(496)} %                        Remove in Eng. 


%\nomer{Vol.\,47, No.\,9}%               Set in Eng. 


%\str{\pageref{begin}--\pageref{end}}%  Set in Eng. 


\udk{517.51} 


 


\author{\it Т.В.\,Ершова} 


% NB:   ^^ remove in Eng. 


\title{Асимптотики для центральных моментов модификаций операторов, 


подобных многочленам Бернштейна} 


\thanks{}


 


\date{} 


%\pagestyle{myheadings}%         Set in Eng. 


\pagestyle{plain} %              Remove in Eng. 


\begin{document} 


%\thispagestyle{plain}%          Set in Eng. 


\maketitle 


\label{begin} 


 


 


 


В данной статье, написанной на материале депонированной работы [1],


модификации B.C.\,Ви\-денского 


и Т.П.\,Пендиной применяются к линейным операторам, 


которые имеют центральные моменты, сходные с центральными моментами 


многочленов Бернштейна. Установлены асимптотики для центральных 


моментов модификаций этих операторов, а также доказаны теоремы типа 


теорем Вороновской--Бернштейна, 


имеющие локальный характер. 


\medskip 


 


{\bf 1.} Как известно ([2], с.\,52), модификации B.C.\,Виденского 


и Т.П.\,Пендиной операторов $L_n$, 


$n\in N$, заданных в пространстве $C[0,1]$, определяются следующими 


рекуррентными формулами: 


\begin{equation} 


\begin{split} 


L_{n1}f&=L_n f, \\ 


L_{n\nu}f&=L_n f-\sum^{\nu-1}_{i=1}\frac{S_i(L_n)}{i!} 


L_{n,\nu-i}f^{(i)},\quad \nu\ge2. 


\end{split} 


\end{equation} 


Предполагаем, что $L_n1=1$ и оператор $L_{n\nu}$ действует в пространстве 


$C^{(\nu-1)}[0,1]$. Функции 


$S_i(L_n,x)=L_n((t-x)^i,x)$, $i=0,1,\dotsc$, называются центральными 


моментами порядка $i$ 


оператора $L_n$. Аналогично определяются центральные моменты 


модификаций $L_{n\nu}$. Полагая 


$f(t)=(t-x)^\nu$, из (1) получаем 


\begin{equation} 


S_\nu(L_{n\nu})=S_\nu(L_n)-\sum^{\nu-1}_{i=1} 


C^i_\nu S_i(L_n)S_{\nu-i}(L_{n,\nu-i}) 


\end{equation} 


или 


\begin{equation} 


\frac{S_\nu(L_{n\nu})}{\nu!}=\frac{S_\nu(L_n)}{\nu!} 


-\sum^{\nu-1}_{i=1}\frac{S_i(L_n)}{i!} 


\frac{S_{\nu-i}(L_{n,\nu-i})}{(\nu-i)!}. 


\end{equation} 


Напомним, что для центральных моментов четных порядков линейных 


положительных операторов (л.\,п.\,о.) имеет место неравенство [3] 


\begin{equation} 


S_{2p}S_{2q}\le S_{2(p+q)},\quad p,q\in N. 


\end{equation} 


 


Условимся считать, что утверждения леммы 1 имеют место в точке $x$ 


из $[0,1]$, и для 


удобства записи символ $x$ писать не будем. 


 


\begin{lem} 


Пусть $\{L_n\}$ --- последовательность л.\,п.\,о., 


$S_{2k-1}(L_n)=O(S_{2k}(L_n))$, $k\ge1$. 


Тогда $S_{2k-1}(L_{n,2k-1})=O(S_{2k}(L_n))$ и 


$S_{2k}(L_{n,2k})=O(S_{2k}(L_n))$. 


\end{lem} 


 


\begin{pf} 


Из (2) при $\nu=2k-1$ и $\nu=2k$ следует соответственно 


\begin{gather} 


S_{2k-1}(L_{n,2k-1})=S_{2k-1}(L_n)-\sum^{2k-2}_{i=1} 


C^i_{2k-1}S_i(L_n)S_{2k-1-i}(L_{n,2k-1-i}), \\ 


% 


S_{2k}(L_{n,2k})=S_{2k}(L_n)-\sum^{2k-1}_{i=1} 


C^i_{2k}S_i(L_n)S_{2k-i}(L_{n,2k-i}). 


\end{gather} 


Доказательство проведем по индукции. Пусть $k=1$. Тогда 


$S_1(L_{n1})=S_1(L_n)=O(S_2(L_n))$. 


Для центрального момента $S_2(L_{n2})$ согласно (6) имеем 


$S_2(L_{n2})=S_2(L_n)-2S^2_1(L_n)$. Отсюда 


$S_2(L_{n2})=S_2(L_n)-O(S_2(L_n)) O(S_2(L_n))= 


S_2(L_n)-O(S_2(L_n))=O(S_2(L_n))$. Предположим, что лемма 


доказана для $k\le m$. Пусть $k=m+1$. Равенство (5) примет вид 


\begin{equation} 


S_{2m+1}(L_{n,2m+1})=S_{2m+1}(L_n)-\sum^{2m}_{i=1} 


C^i_{2m+1}S_i(L_n)S_{2m+1-i}(L_{n,2m+1-i}). 


\end{equation} 


Рассмотрим отдельно случаи четного и нечетного $i$. 


Если $i=2j$, то, в силу (4)


$$ 


S_{2j}(L_n)S_{2m+1-2j}(L_{n,2m+1-2j})=S_{2j}(L_n) 


O(S_{2m-2j+2}(L_n))=O(S_{2j}(L_n)S_{2m-2j+2}(L_n))= 


O(S_{2m+2}(L_n)). 


$$ 


Если $i=2j-1$, то с учетом (4)


\begin{multline*} 


S_{2j-1}(L_n)S_{2m-2j+2}(L_{n,2m-2j+2}) 


=O(S_{2j}(L_n))S_{2m-2j+2}(L_{n,2m-2j+2})= \\ 


% 


=O(S_{2j}(L_n))O(S_{2m-2j+2}(L_n))=O(S_{2m+2}(L_n)). 


\end{multline*} 


Таким образом, из (7) следует $S_{2m+1}(L_{n,2m+1})=O(S_{2m+2}(L_n))$. 


Запишем равенство (6) при 


$k=m+1$: 


\begin{equation} 


S_{2m+2}(L_{n,2m+2})=S_{2m+2}(L_n)-\sum^{2m+1}_{i=1} 


C^i_{2m+2}S_i(L_n)S_{2m+2-i}(L_{n,2m+2-i}). 


\end{equation} 


Если $i=2j$, то 


$$ 


S_{2j}(L_n)S_{2m+2-2j}(L_{n,2m+2-2j})=S_{2j}(L_n) 


O(S_{2m+2-2j}(L_n))=O(S_{2m+2}(L_n)). 


$$ 


Для нечетного $i=2j-1$ имеем 


$$ 


S_{2j-1}(L_n)S_{2m-2j+3}(L_{n,2m-2j+3})=O(S_{2j}(L_n) 


S_{2m-2j+4}(L_n))=O(S_{2m+4}(L_n)). 


$$ 


Следовательно, из (8) получаем $S_{2m+2}(L_{n,2m+2})=O(S_{2m+2}(L_n))$. 


\end{pf} 


 


\begin{note} 


Отметим, что в сумме из правой части равенства (6) слагаемые с нечетным 


индексом $i$ имеют, вообще говоря, порядок малости выше, чем 


слагаемые с четным индексом $i$. 


\end{note} 


 


\begin{note} 


Условиям леммы удовлетворяют многочлены Бернштейна, Канторовича [2] 


и многие другие полиномиальные операторы [4]. 


\end{note} 


 


 


{\bf 2.} Рассмотрим частные суммы $\sigma_i$ ряда Маклорена функции 


$y=e^x$ при $x=a$, где $a$ --- 


произвольное число: $\sigma_i=\sum\limits^i_{k=0}\frac{a^k}{k!}$. Будем 


считать, что $\sigma_i=0$, если $i<0$. 


 


\begin{lem} 


Имеют место равенства 


\begin{alignat}{2} 


\sum^{2i-1}_{j=0}(-1)^j\sigma_j\sigma_{2i-1-j}&=0, &\quad i&\ge 1, \\ 


\sum^{2i}_{j=0}(-1)^j\sigma_j\sigma_{2i-j}&=1, &\quad i&\ge 0. 


\end{alignat} 


\end{lem} 


 


\begin{pf} 


Обозначим суммы из правых частей равенств (9) и (10) через $A$ и $B$. 


Тогда 


\begin{multline*} 


A=\sum^{i-1}_{j=0}(-1)^j\sigma_j\sigma_{2i-1-j} 


+\sum^{2i-1}_{k=i}(-1)^k\sigma_k\sigma_{2i-1-k}\ (k=2i-1-j)= \\ 


% 


=\sum^{i-1}_{j=0}(-1)^j\sigma_j\sigma_{2i-1-j}- 


\sum^0_{j=i-1}(-1)^j\sigma_{2i-1-j}\sigma_j=0. 


\end{multline*} 


В [5] доказано равенство $\sum\limits^{2i}_{j=0}(-1)^j 


(\sigma_j-\sigma_{j-2})\sigma_{2i-j}=0$, что равносильно 


$\sum\limits^{2i}_{j=0}(-1)^j\sigma_j\sigma_{2i-j}= 


\sum\limits^{2i-2}_{j=0}(-1)^j\sigma_j\sigma_{2i-2-j}$. 


Следовательно, $B=\sum\limits^0_{j=0}(-1)^j\sigma_j\sigma_{-j}= 


\sigma^2_0=1$. 


\end{pf} 


 


 


В лемме 3 установим асимптотику для 


центральных моментов порядка $2m$ модификаций 


$L_{n,2m}$ операторов $L_n$, 


центральные моменты которых удовлетворяют некоторым условиям. 


 


\begin{lem} 


Пусть $\{L_n\}$ --- последовательность 


л.\,п.\,о.$;$ $S_{2k-1}(L_n,x)=O(S_{2k}(L_n,x));$ 


\begin{equation} 


\frac{S_{2k}(L_n,x)}{(2k)!}=\sum^k_{i=0} 


A_i\sigma_{k-i}g^k(n)\varphi^k(x)+o(g^k(n)), 


\end{equation} 


$A_i\in\Bbb R$, $A_0=1$, $k\ge1$, $g(n)$ и $\varphi(x)$ --- функции, 


причем $\lim\limits_{n\to\infty}g(n)=0$. Тогда 


\begin{equation} 


\frac{S_{2k}(L_{n,2k},x)}{(2k)!}=(-1)^{k+1}\sum^k_{i=0} 


\alpha_i\sigma_{k-i}g^k(n)\varphi^k(x)+o(g^k(n)), 


\end{equation} 


$\alpha_0=1$, $\alpha_1=A_1\alpha_0$, 


$\alpha_2=A_1\alpha_1-A_2\alpha_0-1$, 


$\alpha_i=\sum\limits^i_{j=1}(-1)^{j+1}A_j\alpha_{i-j}$, $i\ge3$. 


\end{lem} 


 


 


{\bf Доказательство.} Применим индукцию по $k$. В силу условий леммы и 


формулы (3) 


при $\nu=2$ соотношение (12)  имеет место при $k=1$. Предположим, что для 


$k\le m-1$ 


соотношение (12) выполняется. 


Как и в лемме 1, рассмотрим отдельно 


четные и нечетные слагаемые из суммы в правой части (3) для $\nu=2m$. 


Полагая $i=2j$, $j=1,2,\dotsc,m-1$, и применяя (11) 


и предположение индукции, для соответствующего слагаемого имеем 


$$


\frac{S_{2j}(L_n,x)}{(2j)!} 


\frac{S_{2m-2j}(L_{n,2m-2j},x)}{(2m-2j)!}= 


% 


(-1)^{m-j+1}\sum^j_{i=0}A_i\sigma_{j-i} 


\sum^m_{k=0}\alpha_k\sigma_{m-j-k} 


g^m(n)\varphi^m(x)+o(g^m(n)). 


$$


Для нечетных $i=2j-1$ получим 


$$ 


S_{2j-1}(L_n,x)S_{2m-2j+1}(L_{n,2m-2j+1},x)= 


O(S_{2j}S_{2m-2j+2})=O(g^{m+1}(n))=o(g^m(n)). 


$$ 


Следовательно, после очевидных преобразований 


соотношения (3) при $\nu=2m$ приходим к задаче доказательства равенства 


\begin{equation} 


\sum\equiv\sum^m_{j=0}(-1)^j\sum^j_{i=0}A_i\sigma_{j-i} 


\sum^m_{k=0}\alpha_k\sigma_{m-j-k}=0. 


\end{equation} 


Изменив в сумме (13) порядок суммирования, получим 


$$ 


\sum=\sum^m_{k=0}\alpha_k\sum^m_{i=0}A_i\sum^m_{j=i} 


(-1)^j\sigma_{j-i}\sigma_{m-k-j}. 


$$ 


Внутренняя сумма после перехода к переменной суммирования $t=j-i$ будет 


равна 


$$ 


(-1)^i\sum^{m-i}_{t=0}(-1)^t\sigma_t\sigma_{m-i-k-t}=(-1)^i 


\sum^{m-i-k}_{t=0}(-1)^t\sigma_t\sigma_{m-i-k-t}, 


$$ 


т.\,к. $\sigma_p=0$ при $-k\le p\le-1$. Применяя лемму 2, получим 


$\sum=\sum\limits^m_{i,k=0}{\!\!\!'}(-1)^i A_i\alpha_k$, где \ ${}'$ \


означает, что суммирование производится по $i$ и $k$, для которых $m-i-k$ 


четно и неотрицательно. Покажем, что для любого $m$ сумма $\sum$ 


равна нулю. Действительно, если $m=1$, то 


$\sum=A_0\alpha_1-A_1\alpha_0=0$. Если $m=2$, то 


$\sum=A_0\alpha_0+A_0\alpha_2-A_1\alpha_1 


+A_2\alpha_0=1+\alpha_2-A_1\alpha_1+A_2\alpha_2=0$. 


Для $m\ge3$ имеем 


$$ 


\sum=\sideset{}{'}\sum^m_{i,k=0}(-1)^i A_i\alpha_k= 


\sideset{}{'}\sum^{m-2}_{i,k=0}(-1)^i A_i\alpha_k+ 


\sum_{i+k=m}(-1)^i A_i\alpha_k=\sum^m_{i=0}(-1)^i 


A_i\alpha_{m-i}=0.\quad\square 


$$ 


 


\begin{note} 


Из леммы 3 следуют результаты предыдущих работ автора. Для многочленов 


Бернштейна и Канторовича имеет место равенство [4], [6] 


$\frac{S_{2k}(L_n,x)}{(2k)!}= \frac{x^k(1-x)^k}{2^k k!n^k}+ 


o(\frac{1}{n^k})$. 


Соотношение (12) вида $\frac{S_{2k}(L_{n,2k},x)}{(2k)!}= (-1)^{k+1} 


\frac{x^k(1-x)^k}{2^k k!n^k}+o(\frac{1}{n^k})$ доказано в [7]. 


Многочлены $U_n$, введенные в ([8], с.\,48), 


удовлетворяют условию (11): $\frac{S_{2k}(U_n,x)}{(2k)!}= (\frac{1}{2^k 


k!}+ \frac{1}{2^{k-1}(k-1)!}) \frac{x^k(1-x)^k}{n^k}$. 


Для них соотношение $\frac{S_{2k}(U_{n,2k},x)}{(2k)!}= (-1)^{k+1}\sigma_k 


\frac{x^k(1-x)^k}{n^k}$ доказано в [5]. 


В случае этих трех операторов 


лемма 3 имеет место в интервале $(0,1)$ с числом $a=1/2$. 


\end{note} 


 


\begin{lem} 


Пусть $\{L_n\}$ --- последовательность л.\,п.\,о., удовлетворяющих 


условиям $(11)$ и 


\begin{equation} 


\frac{S_{2k-1}(L_n,x)}{(2k-1)!}=\sum^k_{i=1}B_i\sigma_{k-i} 


g^k(n)\varphi^k(x)h(x)+o(g^k(n)), 


\end{equation} 


$B_i\in\Bbb R$, $k\ge1$, $h(x)$ --- функция. Тогда 


\begin{gather} 


\frac{S_{2k-1}(L_{n,2k-1},x)}{(2k-1)!}=(-1)^{k+1} 


\sum^k_{i=1}\beta_i\sigma_{k-i}g^k(n) 


\varphi^k(x)h(x)+o(g^k(n)), \\ 


% 


\beta_i=\sum^{i-1}_{j=1}(-1)^{j+1}A_j \beta_{i-j} 


+\sum^i_{j=1}(-1)^{j+1}B_j\alpha_{i-j}.\notag 


\end{gather} 


\end{lem} 


 


\begin{pf} 


Применим индукцию по $k$. При $k=1$ (15) совпадает с (14), т.\,к. 


$\beta_1=B_1$. Пусть для $k\le m-1$ имеет место соотношение (15). 


Рассмотрим $k=m$. Положим 


в формуле (3) $\nu=2m-1$. Тогда 


\begin{equation} 


\frac{S_{2m-1}(L_{n,2m-1})}{(2m-1)!}= 


\frac{S_{2m-1}(L_n)}{(2m-1)!}-\sum^{2m-2}_{i=1} 


\frac{S_i(L_n)}{i!} 


\frac{S_{2m-1-i}(L_{n,2m-1-i})}{(2m-1-i)!}. 


\end{equation} 


Полагая $i=2j$ и $i=2j-1$, получим $j=1,\dotsc,m-1$. Выполним 


преобразования 


равенства (16), используя (11), (12), (14) 


и предположение индукции. 


Остается установить тождество 


$$ 


\sum_1+\sum_2\equiv\sum^{m-1}_{j=0}(-1)^j\sum^j_{p=0} 


A_p\sigma_{j-p}\sum^{m-j}_{t=1}\beta_t\sigma_{m-j-t} 


+\sum^m_{j=1}(-1)^j\sum^j_{l=1}B_l\sigma_{j-l} 


\sum^{m-j}_{s=0}\alpha_s\sigma_{m-j-s}=0. 


$$ 


Слагаемые, стоящие в скобках первой суммы 


$$ 


\sum_1=\sum^{m-1}_{j=0}(-1)^j\bigg(\,\sum^j_{p=0} 


\sum^{m-j}_{t=1}A_p\beta_t\sigma_{j-p}\sigma_{m-j-t}\bigg), 


$$ 


объединим в группы в соответствии со значением $q=p+t$. 


Замечая, что $p+t=1,\dotsc,m$, имеем 


$$ 


\sum_1=\sum^{m-1}_{j=0}(-1)^j\sum^m_{q=1}\sum^{q-1}_{p=0} 


A_p\beta_{q-p}\sigma_{j-p}\sigma_{m-j-q+p}= 


\sum^m_{q=1}\sum^{q-1}_{p=0}A_p\beta_{q-p}\sum^{m-1}_{j=0} 


(-1)^j\sigma_{j-p}\sigma_{m-j-q+p}. 


$$ 


Для вычисления последней внутренней суммы $\sum\limits_3$ используем лемму 


2. Прежде всего заметим, 


что $\sum\limits_3=\sum\limits^{m-q+p}_{j=p}(-1)^j 


\sigma_{j-p}\sigma_{m-j-q+p}$. Полагая $t=j-p$, получим 


$\sum\limits_3=(-1)^p \sum\limits^{m-q}_{t=0}(-1)^t \sigma_t 


\sigma_{m-q-t}$, и 


$\sum\limits_3$ согласно лемме 2 отлична от нуля при условии, что $m$ и 


$q$ 


одинаковой четности. В этом случае 


$\sum\limits_3=(-1)^p$. Таким образом, 


$$ 


\sum_1=\sideset{}{'}\sum^m_{q=1}\sum^{q-1}_{p=0} 


(-1)^pA_p\beta_{q-p}, 


$$ 


где \ ${}'$ \ означает, что суммирование 


производится по тем $q$, четность которых совпадает 


с четностью $m$. Аналогично устанавливаем, что 


$$ 


\sum_2=\sideset{}{'}\sum^m_{q=1}\sum^{q}_{l=1} 


(-1)^l B_l\alpha_{q-l}. 


$$ 


Здесь \ ${}'$ \ также означает суммирование по $q$, имеющим 


одинаковую четность с $m$. Следовательно, 


\begin{equation} 


\sum_1+\sum_2=\sideset{}{'}\sum^m_{q=1}\bigg(\, 


\sum^{q-1}_{p=0}(-1)^p A_p\beta_{q-p}+ 


\sum^q_{l=1}(-1)^l B_l\alpha_{q-l}\bigg). 


\end{equation} 


Согласно определению $\beta_i$ сумма в скобках, 


стоящая в правой части (17), равна нулю. 


\end{pf} 


 


\begin{note} 


Лемма 4 для многочленов Бернштейна доказана в [9], для многочленов 


Канторовича --- в [10] при любом $x$ из $(0,1)$. 


\end{note} 


 


\begin{defnnonum} 


Операторы $L_n$, удовлетворяющие условиям (11) и (14) на $(0,1)$, назовем 


операторами, подобными многочленам Бернштейна. 


\end{defnnonum} 


 


{\bf 3.} 


{\bf Лемма 5.} {\it


Пусть последовательность л.\,п.\,о. $\{L_n\}$ удовлетворяет условию 


$(11)$. Тогда}


\begin{equation} 


S^*_{2k-1}(L_n,x)\asymp g^{k-\frac{1}{2}}(n). 


\end{equation} 


\addtocounter{lem}{1}


 


\begin{pf} 


Будем использовать неравенство Г\"ельдера для л.\,п.\,о. 


\begin{equation} 


L_n(u v)\le(L_n u^p)^{\frac{1}{p}} 


(L_n v^q)^{\frac{1}{q}}, 


\end{equation} 


где функции $u$, $v$ неотрицательны, $p>1$, $q>1$, 


$\frac{1}{p}+\frac{1}{q}=1$. Из этого неравенства легко 


получить следующие неравенства: 


\begin{align} 


S^*_{2k-1}(L_n)&\le\sqrt{S_{2k-2}(L_n)S_{2k}(L_n)}, \\ 


S^3_{2k-2}(L_n)&\le (S^*_{2k-1}(L_n))^2 


S_{2k-4}(L_n). 


\end{align} 


В первом случае в (19) надо положить $p=q=2$, $u(t)=|t-x|^{k-1}$, 


$v(t)=|t-x|^k$, 


во втором положим $p=3/2$, $q=3$, $u(t)=|t-x|^{2/3(2k-1)}$, 


$v(t)=|t-x|^{\frac{2k-4}{3}}$. 


В силу (11) из (20) и (21) 


следуют неравенства 


$$ 


S^*_{2k-1}(L_n,x)\le C_1 g^{k-\frac{1}{2}}(n)\ \ 


\text{и}\ \ S^*_{2k-1}(L_n,x)\ge C_2 


g^{k-\frac{1}{2}}(n), 


$$ 


где $C_1$ и $C_2$ --- некоторые постоянные, не зависящие от $n$. 


\end{pf} 


 


\begin{lem} 


Пусть для центральных моментов л.\,п.\,о. $L_n$ 


выполняется условие $(11)$ и 


функция $f\in C^{(\nu-1)}[0,1]$ дифференцируема в точке $x$ \ $\nu+l$ 


$(\nu\ge1$, $l\ge-1)$ раз. Тогда 


\begin{equation} 


L_{n\nu}(f,x)=f(x)+\sum^{\nu+l}_{i=\nu} 


\frac{f^{(i)}(x)}{i!}S_i(L_{n\nu},x)+o(g^{\frac{\nu+l}{2}}(n)). 


\end{equation} 


Сумму из правой части $(22)$ при $l=-1$ считаем отсутствующей. 


\end{lem} 


 


\begin{pf} 


В силу (11) и (18) выполняется $S^*_i(L_n,x)\asymp g^{i/2}(n)$, что позволяет 


применить теорему 2 из [11] к функции $f$ и модификациям $L_{n\nu}$. 


Таким образом, имеем 


$$ 


L_{n\nu}(f,x)=f(x)+\sum^{\nu+l}_{i=\nu}\frac{f^{(i)}(x)}{i!} 


S_i(L_{n\nu},x)+o(S^*_{\nu+l}(L_n)). 


$$ 


В последнем слагаемом осуществляем соответствующую замену. 


\end{pf} 


 


 


Следующие две теоремы являются асимптотическими 


теоремами Вороновской--Бернштейна для исследуемых в данной статье 


модификаций линейных операторов $L_{n,2m}$ и $L_{n,2m-1}$. 


 


\begin{thm} 


Пусть для последовательности л.\,п.\,о. $\{L_n\}$ 


выполнены условия леммы $3$ и 


функция $f\in C^{(2m-1)}[0,1]$, $m\ge1$, в точке $x$ дифференцируема $2m$ 


раз. Тогда 


$$ 


L_{n,2m}(f,x)=f(x)+f^{(2m)}(x)(-1)^{m+1}\sum^m_{i=0} 


\alpha_i\sigma_{m-i}g^m(n)\varphi^m(x)+o(g^m(n)). 


$$ 


\end{thm} 


 


\begin{pf} 


В соотношении (22) положим $\nu=2m$, $l=0$, а затем применим (12). 


\end{pf} 


 


\begin{thm} 


Пусть для последовательности л.\,п.\,о. $\{L_n\}$ выполнены условия леммы 


$4$, 


функция $f\in C^{(2m-2)}[0,1]$, $m\ge1$, в точке $x$ дифференцируема $2m$ 


раз. Тогда 


\begin{multline*} 


L_{n,2m-1}(f,x)=f(x)+f^{(2m-1)}(x)(-1)^{m+1} 


\sum^m_{i=1}\beta_i\sigma_{m-i}g^m(n)\varphi^m(x)h(x)+ \\ 


% 


+ 


f^{(2m)}(x)(-1)^{m+1}\sum^m_{i=0}\alpha_i\sigma_{m-i} 


g^m(n)\varphi^m(x)+o(g^m(n)). 


\end{multline*} 


\end{thm} 


 


\begin{pf} 


Требуемое соотношение следует из (22) при $\nu=2m-1$, $l=1$ и (12), 


(15) при $k=m$. 


\end{pf} 


 


 


\begin{thebibliography}{99} 


 


\bibitem{1} 


Ершова Т.В. {\em Асимптотики для центральных моментов модификаций 


операторов, 


подобных многочленам Бернштейна}. -- Челяб. гос. пед. ун-т. -- 


Челябинск, 2000. -- 14 с. -- Деп. 


в ВИНИТИ 03.05.00, \No\,1277-B00. 


\bibitem{2} 


Виденский B.C. {\em Многочлены Бернштейна}. -- Л.: Изд-во 


Ленингр. гос. пед. ин-та, 1990. -- 63~с. 


\bibitem{3} 


Ершова Т.В. {\em Об итерациях линейных положительных 


операторов}. -- Челяб. гос. пед. 


ун-т. -- Челябинск, 1999. -- 10 с. -- Деп. в ВИНИТИ 10.03.99, 


\No\,786-B99. 


\bibitem{4} 


Волков Ю.И. {\em О некоторых линейных положительных операторах} 


// Матем. заметки. -- 


1978. -- Т.\,237. -- \No\,5. -- С.\,659--669. 


\bibitem{5} 


Ершова Т.В. {\em О модификациях многочленов, подобных многочленам 


Бернштейна} // Вестн. 


Челяб. ун-та. Сер.\,3. -- 1996. -- \No\,1. -- С.\,49--54. 


\bibitem{6} 


Бернштейн С.Н. {\em Добавление к статье Е.В.\,Вороновской 


``Определение асимптотического вида приближения функций 


полиномами С.Н.\,Бернштейна''}. -- M.: Изд-во 


АН СССР, 1954. -- Т.\,2. -- С.\,155--158. 


\bibitem{7} 


Ершова Т.В. {\em О некоторых свойствах модификаций многочленов Бернштейна} 


// Функциональный анализ. -- Ульяновск, 1993. -- \No\,34. -- С.\,15--26. 


\bibitem{8} 


Виденский B.C. {\em Линейные положительные операторы 


конечного ранга}. -- Л.: Изд-во 


Ле\-нингр. гос. пед. ун-та, 1985. -- 68 с. 


\bibitem{9} 


Ершова Т.В. {\em Асимптотическая теорема Вороновской--Бернштейна 


для модификаций 


многочленов Бернштейна} // Вестн. Челяб. yн-тa. Сер.\,3. -- 


1994. -- \No\,1. -- С.\,43--48. 


\bibitem{10} 


Ершова Т.В. {\em Предельные соотношения модификаций многочленов 


Канторовича}. -- Челяб. гос. пед. ун-т. -- Челябинск, 1996. -- 


9 с. -- Деп. в ВИНИТИ 28.10.96, \No\,3146-B96. 


\bibitem{11} 


Ершова Т.В. {\em О приближении непрерывных функций модификациями линейных 


положительных операторов} // Прим. функционального анализа 


в теории приближений. -- Тверь, 


1997. -- С.\,73--79. 


 


\end{thebibliography} 


 


\vskip 2cm 


 


\postup{\it Челябинский государственный}{\qquad \it Поступили} 


\postup{\it педагогический университет}{\it первый вариант $10.05.2001$} 


\postup{}{\it окончательный вариант $07.04.2003$} 


 


\label{end} 


\end{document} 


